\documentclass[12pt,letterpaper]{article}

%% ── Packages ──────────────────────────────────────────────────────────────────
\usepackage[margin=1in]{geometry}
\usepackage{amsmath,amssymb}
\usepackage{booktabs}
\usepackage{graphicx}
\usepackage[numbers,sort&compress]{natbib}
\usepackage{setspace}
\usepackage{microtype}
\usepackage{xcolor}
\usepackage{siunitx}
\usepackage{array}
\usepackage{longtable}
\usepackage{caption}
\usepackage{authblk}
\usepackage{lineno}
\usepackage{tabularx}
\usepackage{hyperref}

%% ── Hyperref setup ────────────────────────────────────────────────────────────
\hypersetup{
  colorlinks = true,
  linkcolor  = blue,
  citecolor  = blue,
  urlcolor   = blue
}

%% ── Line numbers ──────────────────────────────────────────────────────────────
\linenumbers

%% ── Spacing ───────────────────────────────────────────────────────────────────
\doublespacing

%% ── Custom commands ───────────────────────────────────────────────────────────
%% All custom shorthands wrap content in $...$ so \mathrm works everywhere
\newcommand{\SCE}{$\mathrm{SCE}$}
\newcommand{\SCEstar}{$\mathrm{SCE}^{*}$}
\newcommand{\CERT}{$\mathrm{CE}_{\mathrm{RT}}$}
\newcommand{\CELT}{$\mathrm{CE}_{\mathrm{LT}}$}
\newcommand{\Rsq}{$R^{2}$}
\newcommand{\dSCE}{$\Delta\mathrm{SCE}$}
%% Math-mode only versions for use inside equations
\newcommand{\mSCE}{\mathrm{SCE}}
\newcommand{\mSCEstar}{\mathrm{SCE}^{*}}
\newcommand{\mCERT}{\mathrm{CE}_{\mathrm{RT}}}
\newcommand{\mCELT}{\mathrm{CE}_{\mathrm{LT}}}

%% ════════════════════���══════════════════════════════════════════════════════════
\begin{document}

\title{\textbf{Solvation Configuration Entropy Predicts
Lithium Metal Battery Coulombic Efficiency:\\
A Mathematically Derived Optimal Coordination Disorder
and the Dual-Temperature Performance Band}}

\author[1]{Eric Robert Lawson}
\affil[1]{OrganismCore, Independent Research.
  ORCID: 0009-0002-0414-6544.
  Pre-registration timestamp: 2026-04-01,
  \texttt{github.com/Eric-Robert-Lawson/attractor-oncology}}

\date{Version 1.0 \quad \today}

\maketitle
\thispagestyle{empty}

%% ── Abstract ────────────────────────────���─────────────────────────────────────
\begin{abstract}
We introduce the \textbf{Solvation Configuration Entropy}
($\mSCE$) --- the Shannon entropy of the
$\mathrm{Li}^{+}$ first-shell coordination geometry
population distribution --- as a quantitative predictor
of lithium metal battery Coulombic efficiency (CE) across
both room temperature (RT) and low temperature
(LT, $\leq -20\,^{\circ}\mathrm{C}$).
Applying the framework to 21 published electrolyte systems,
we confirm four empirical equations with statistical
significance ($R^{2} = 0.828$,
$p = 4.5 \times 10^{-6}$, $n = 17$
for the two-variable model;
$R^{2} = 0.708$, $p = 0.009$
for the Regime~1 log-deficit equation).
From these confirmed equations we derive analytically that a
\textbf{mathematically optimal} $\mSCEstar = 1.466$ exists,
constituting a global maximum of combined RT$+$LT performance
(second derivative always negative).
At $\mSCEstar$, the required solvation structure is
$\mathrm{dom\%} \approx 38\%$ with three significant
coordination configurations.
We identify a \textbf{dual-temperature performance band}
($\mSCE \approx 1.45$--$1.47$) within which electrolytes
simultaneously achieve
$\mathrm{CE}_{\mathrm{RT}} \geq 90\%$ and
$\mathrm{CE}_{\mathrm{LT}} \geq 60\%$
at $-20\,^{\circ}\mathrm{C}$:
a rare outcome confirmed by a systematic literature survey
as present in fewer than five published systems,
none designed toward the band from first principles.
A literature search of the 2018--2026 corpus using
SES~AI's Molecular Universe platform confirms that no
published work has defined $\mSCE$, correlated it with CE,
or derived a mathematical performance optimum from it.
We further predict that the field's two best
room-temperature performers --- HFTHP and BTFMD ---
achieve near-zero $\mSCE$ and will show
$\mathrm{CE}_{\mathrm{LT}} \approx 22\%$
at $-20\,^{\circ}\mathrm{C}$,
a measurement absent from the published literature.
Two novel candidate electrolytes designed toward
$\mSCEstar$ are proposed with specific,
falsifiable performance predictions.
\end{abstract}

\newpage
\tableofcontents
\newpage

%% ═══════════════════════════════════════════════════════════════════════════════
\section{Introduction}
\label{sec:intro}
%% ═══════════════════════════════════════════════════════════════════════════════

Lithium metal batteries offer the highest theoretical energy
density of any anode chemistry, yet their practical deployment
is constrained by two competing failure modes that act across
temperature.
At room temperature, dendrite nucleation and incomplete
solid-electrolyte interphase (SEI) formation limit
Coulombic efficiency and cycle life.
At low temperature ($\leq -20\,^{\circ}\mathrm{C}$),
sluggish $\mathrm{Li}^{+}$ desolvation kinetics cause
capacity fade and resistance rise even in electrolytes
that perform well at room temperature~\cite{holoubek2022,fan2023}.

The dual-temperature performance problem --- achieving
$\mathrm{CE}_{\mathrm{RT}} \geq 90\%$ and
$\mathrm{CE}_{\mathrm{LT}} \geq 60\%$ simultaneously
--- is rare in the published literature.
A systematic survey of the 2022--2025 record found fewer
than five systems with quantitative CE data at both
temperatures, none designed toward an explicit
dual-temperature target~\cite{hu2025,zhai2025,
natcomm2025motif,natenergy2023hee}.

The dominant approach to electrolyte design has been
correlation-guided: screen candidate molecules for intrinsic
properties (HOMO/LUMO energies, donor number, solvation
energy), synthesise, test, and iterate.
This approach has produced excellent room-temperature
electrolytes, including HFTHP- and BTFMD-based systems
with $\mathrm{CE}_{\mathrm{RT}} > 99\%$~\cite{lee2024btfmd}.
It has not produced a principled design target for combined
RT$+$LT performance, because the property governing that
combination --- the \textit{statistical diversity of the
$\mathrm{Li}^{+}$ coordination geometry population} ---
has not been defined as a quantitative descriptor.

Recent work has approached this space from adjacent
directions.
Hu et al.~\cite{hu2025} use entropy qualitatively to
describe compositional diversity in multi-component
electrolytes.
Zhai et al.~\cite{zhai2025} invoke a ``high-entropy
effect'' to describe weakened coordination.
Most directly, a concurrent Nature Communications
paper~\cite{natcomm2025motif} introduced
\textit{interaction motif descriptors} --- a systematic
framework for categorising $\mathrm{Li}^{+}$ coordination
interaction types --- and used them to rationally design
an electrolyte achieving 99.7\% RT CE with documented LT
performance.
Their framework identifies the dominant coordination motif
associated with optimal room-temperature performance.
None of these approaches computes the Shannon entropy of
the full coordination microstate distribution, correlates
it with CE, or derives an optimal entropy value.
A literature search using SES~AI's Molecular Universe
platform confirms that no published paper has done so
across the entire 2018--2026
record~\cite{mu2026query5}.

The Joule 2025 paper from Hunan
University~\cite{joule2025hunan} independently defined a
variable called $S_{\mathrm{sc}}$ (solvation
configurational entropy) and showed that higher
$S_{\mathrm{sc}}$ correlates with improved low-temperature
interfacial kinetics.
They did not derive an optimal value, did not model
room-temperature performance as a function of the same
variable, and did not identify a band.

Here we present the \textbf{Solvation Configuration
Entropy ($\mSCE$)} framework, which:

\begin{enumerate}
  \item Defines $\mSCE$ formally as the Shannon entropy
        of the $\mathrm{Li}^{+}$ first-shell coordination
        microstate population distribution.
  \item Confirms four empirical equations across 21
        published electrolyte systems relating $\mSCE$
        to RT and LT CE.
  \item Derives analytically from those equations that
        $\mSCEstar = 1.466$ is the global maximum of
        combined RT$+$LT performance.
  \item Identifies a dual-temperature performance band
        of width $\approx 0.02$--$0.05$ $\mSCE$ units
        centred near $\mSCEstar$.
  \item Generates seven specific, falsifiable forward
        predictions from the confirmed equations,
        including two novel candidate electrolytes and
        an unmeasured LT CE prediction for HFTHP
        and BTFMD.
\end{enumerate}

%% ═══════════════════════════════════════════════════════════════════════════════
\section{The SCE Variable}
\label{sec:variable}
%% ═══════════════════════════════════════════════════════════════════════════════

\subsection{Definition}
\label{sec:definition}

For a given electrolyte system at fixed composition and
temperature, the $\mathrm{Li}^{+}$ first solvation shell
is occupied by a population of discrete configurations,
each characterised by its coordination numbers
$(n_{\mathrm{solvent}}, n_{\mathrm{anion}})$.
Let $p_{i}$ denote the fraction of $\mathrm{Li}^{+}$
ions occupying configuration $i$ at thermodynamic
equilibrium.
The \textbf{Solvation Configuration Entropy} is:

\begin{equation}
  \mSCE = -\sum_{i} p_{i} \ln p_{i}
  \label{eq:sce}
\end{equation}

$\mSCE$ is zero when all $\mathrm{Li}^{+}$ ions occupy
a single configuration ($p_{1} = 1$), and increases
monotonically as the population spreads across more
configurations of more equal weight.

\subsection{Relationship to SSIP/CIP/AGG}
\label{sec:ssip}

The conventional SSIP/CIP/AGG classification is a
compressed three-bucket representation of the full
coordination population.
$\mSCE$ computed from three buckets alone underestimates
the true $\mSCE$ by the within-bucket entropy
contribution:

\begin{equation}
  \mSCE_{\mathrm{full}} =
  \mSCE_{\mathrm{SSIP/CIP/AGG}}
  + \mSCE_{\mathrm{within\text{-}category}}
  \label{eq:correction}
\end{equation}

For ether-based LiFSI electrolytes, the within-category
correction is approximately 0.3--0.6 $\mSCE$ units,
reflecting the multiple distinct
$(n_{\mathrm{solvent}}, n_{\mathrm{anion}})$
configurations collapsed into each bucket.
All $\mSCE$ values in this paper are computed from the
full configuration population, not from the three-bucket
summary.

\subsection{Step 1 Error and Correction}
\label{sec:step1}

An initial calculation (Step~1) computed Shannon entropy
across chemically heterogeneous RDF pair coordination
numbers.
This is incorrect: it measures chemical complexity of the
RDF file, not solvation geometry variance.
The correct calculation (Step~2 onward) uses the
population distribution of discrete
$(n_{\mathrm{solvent}}, n_{\mathrm{anion}})$
configurations.
The distinction raised $R^{2}$ from 0.336 to 0.815 on
the salt-only dataset ($n = 9$, $p = 0.0009$).

%% ═══════════════════════════════════════════════════════════════════════════════
\section{Dataset and Empirical Analysis}
\label{sec:dataset}
%% ═══════════════════════════════════════════════════════════════════════════════

\subsection{Dataset Construction}
\label{sec:data}

The analysis uses 21 electrolyte systems drawn from the
published literature, with $\mSCE$ values, Coulombic
efficiency data, and solvation structure
characterisation.
Nine systems have $\mSCE$ computed from explicit
configuration population tables extracted from
supplementary data of Energy Advances
2025~\cite{energyadvances2025}.
Twelve systems have $\mSCE$ and CE from the primary
literature~\cite{fan2023,btfmd2022,joule2025hunan,
cao2022,holoubek2022,wan2023,zhang2024,wan2023acsel,
peng2021,mdpibatt2026}.
No estimated values are used in the final analysis.

\subsection{Three-Regime Classification}
\label{sec:regimes}

A critical structural feature of the dataset is that the
relationship between $\mSCE$ and
$\mathrm{CE}_{\mathrm{RT}}$ is not monotonic across the
full dataset.
Three distinct regimes are identified:

\begin{itemize}
  \item \textbf{Regime~1 (R1):}
        $\mathrm{CE}_{\mathrm{RT}} < 90\%$,
        $n = 8$, $\mSCE$ range 1.15--2.09.
        CE is geometry-limited: $\mSCE$ predicts the
        degree to which solvation geometry incompatibility
        degrades SEI completeness.
        Higher $\mSCE \rightarrow$ lower
        $\mathrm{CE}_{\mathrm{RT}}$.

  \item \textbf{Regime~2 (R2):}
        $\mathrm{CE}_{\mathrm{RT}} \geq 90\%$,
        normal LT performance, $n = 9$,
        $\mSCE$ range 1.15--1.74.
        CE has crossed a threshold driven by salt
        concentration and AGG fraction.
        Within R2, higher $\mSCE$ tracks higher AGG
        fraction, inverting the R1 direction.

  \item \textbf{Regime~3 (R3):}
        $\mathrm{CE}_{\mathrm{RT}} \geq 90\%$
        but anomalous LT performance, $n = 3$.
        Either kinetic AGG-locking (BTFMD, TTE) or
        carbonate transport failure
        (LiPF\textsubscript{6}/EC/DMC).
        Excluded from the band analysis.
\end{itemize}

\subsection{Confirmed Equations}
\label{sec:equations}

Eight iterative analytical steps produced four confirmed
equations from the 21-system dataset.

\subsubsection*{Equation 1 --- Regime 1 RT performance}

\begin{equation}
  \ln\!\bigl(100.1 - \mCERT\bigr)
  = 1.493 + 1.230 \times \mSCE
  \label{eq:eq1}
\end{equation}

\noindent Equivalently:
$\mCERT = 100.1 - e^{1.493 + 1.230\,\mSCE}$.

\medskip\noindent
\textit{Statistics:} $R^{2} = 0.708$, $p = 0.009$,
$b_1 = 1.230$, $\mathrm{SE} = 0.322$,
$t(6) = 3.82$, 95\% CI $[0.441, 2.018]$,
Cohen $f^{2} = 2.43$ (large),
$\mathrm{LOO}_{\min}$ $R^{2} = 0.563$,
$n = 8$ (Regime~1 systems).

\subsubsection*{Equation 2 --- Band: LT performance vs SCE}

\begin{equation}
  \mCELT = 11.91 + 33.21 \times \mSCE
  \label{eq:eq2}
\end{equation}

\noindent Valid domain: $\mSCE \approx 1.24$--$2.28$
(normal LT systems, R3 excluded).

\medskip\noindent
\textit{Statistics:} $r = 0.732$, $p = 0.025$,
$n = 9$ (confirmed normal $-20\,^{\circ}\mathrm{C}$
systems + HEE).

\subsubsection*{Equation 3 --- Two-variable model
(full R1 + R2 dataset)}

\begin{equation}
  \mCERT = 104.68
           - 25.85\,\mSCE
           + 31.16\,d_{\mathrm{regime}}
  \label{eq:eq3}
\end{equation}

\noindent where $d_{\mathrm{regime}} = 0$ for Regime~1
and $d_{\mathrm{regime}} = 1$ for Regime~2.

\medskip\noindent
\textit{Statistics:} $R^{2} = 0.828$, $F = 33.65$,
$p = 4.5 \times 10^{-6}$, $n = 17$.

\subsubsection*{Equation 4 --- Within-Regime 2 slope}

\begin{equation}
  \ln\!\bigl(100.1 - \mCERT\bigr)
  = 4.898 - 2.545 \times \mSCE
  \label{eq:eq4}
\end{equation}

\noindent Direction inverted relative to Equation~1:
within R2, higher $\mSCE$ tracks higher AGG fraction,
yielding better RT CE via the concentration-driven
mechanism.

\medskip\noindent
\textit{Statistics:} $R^{2} = 0.474$, $p = 0.040$,
Cohen $f^{2} = 0.900$ (large), $n = 9$.

\subsection{Master Dataset}
\label{sec:masterdata}

\begin{longtable}{llccccc}
\caption{Master gradient table --- 21 electrolyte
systems ordered by $\mSCE$.
RT and LT CE in percent.
LT measured at $-20\,^{\circ}\mathrm{C}$ unless noted.
R = regime. Dom\% = dominant configuration fraction.}
\label{tab:master}\\
\toprule
Rank & System & Conc & R & $\mSCE$ & Dom\% &
$\mathrm{CE}_{\mathrm{RT}}$ \\
\midrule
\endfirsthead
\toprule
Rank & System & Conc & R & $\mSCE$ & Dom\% &
$\mathrm{CE}_{\mathrm{RT}}$ \\
\midrule
\endhead
\bottomrule
\endfoot
1  & FEME/LiFSI           & 4.0\,M & 2   & 1.150 & 47 & 91   \\
2  & LiFSI/DME            & 1.0\,M & 2   & 1.240 & 44 & 97   \\
3  & FEME/LiFSI           & 1.8\,M & 1   & 1.295 & 43 & 82   \\
4  & FEME/LiFSI           & 1.0\,M & 1   & 1.368 & 44 & 70   \\
5  & BTFMD/LiFSI          & 1.0\,M & 3   & 1.401 & 40 & 99.4 \\
6  & LiFSI/DME+FEC        & 1.0\,M & 2   & 1.445 & 38 & 98.0 \\
7  & LiFSI/TTE            & 4.0\,M & 3   & 1.524 & 35 & 99.1 \\
8  & LiFSI/THF            & 1.0\,M & 2   & 1.528 & 30 & 96   \\
9  & LiFSI/2-MeTHF        & 1.0\,M & 2   & 1.552 & 32 & 94.0 \\
10 & LiFSI/DOL            & 1.0\,M & 2   & 1.606 & 30 & 95.8 \\
11 & DPE/LiFSI            & 4.0\,M & 1   & 1.656 & 32 & 75   \\
12 & DPE/LiFSI            & 1.0\,M & 1   & 1.659 & 28 & 55   \\
13 & DPE/LiFSI            & 1.8\,M & 1   & 1.671 & 26 & 65   \\
14 & LiFSI/DME            & 4.0\,M & 2   & 1.703 & 28 & 99.2 \\
15 & LHCE DME/TTE         & 4.0\,M & 2   & 1.735 & 25 & 99.0 \\
16 & LiFSI/DME            & 2.0\,M & 2   & 1.739 & 25 & 98.5 \\
17 & LiPF$_{6}$/EC/DMC    & 1.0\,M & 3   & 1.991 & 20 & 92.0 \\
18 & EC/DEC 1:1           & 1.0\,M & 1   & 2.005 & 24 & 35   \\
19 & EC/DEC 1:1           & 4.0\,M & 1   & 2.010 & 18 & 60   \\
20 & EC/DEC 1:1           & 1.8\,M & 1   & 2.085 & 21 & 40   \\
21 & High-entropy (HEE)   & 1.0\,M & --- & 2.282 & 12 & 93   \\
\end{longtable}

\noindent LT CE values for systems with data:
LiFSI/DME 1.0\,M: 45\%;
LiFSI/DME+FEC: 58\%;
LiFSI/THF: 72\%;
LiFSI/2-MeTHF: 74\%;
LiFSI/DOL: 68\%;
LiFSI/DME 4.0\,M: 78\%;
LHCE: 55\%;
LiFSI/DME 2.0\,M: 62\%;
BTFMD/LiFSI: 30\%;
LiFSI/TTE: 35\%;
LiPF$_{6}$/EC/DMC: 38\%;
HEE: 88\% (at $-40\,^{\circ}\mathrm{C}$).
All others: not published.

%% ═══════════════════════════════════════════════════════════════════════════════
\section{The Mathematically Optimal SCE}
\label{sec:optimum}
%% ═══════════════════════════════════════════════════════════════════════════════

\subsection{Derivation of $\mSCEstar$}
\label{sec:derive}

Define a composite performance score treating RT and LT
performance as equally important:

\begin{equation}
  S(\mSCE) = \mCERT(\mSCE) + \mCELT(\mSCE)
  \label{eq:score}
\end{equation}

Substituting Equations~\ref{eq:eq1} and~\ref{eq:eq2}:

\begin{equation}
  S(\mSCE) = 112.01 + 33.21\,\mSCE
             - e^{1.493 + 1.230\,\mSCE}
  \label{eq:scorefull}
\end{equation}

Taking the first derivative and setting to zero:

\begin{equation}
  \frac{dS}{d\mSCE}
  = 33.21 - 1.230\,e^{1.493 + 1.230\,\mSCE} = 0
  \label{eq:firstderiv}
\end{equation}

\begin{equation}
  e^{1.493 + 1.230\,\mSCE}
  = \frac{33.21}{1.230} = 27.00
  \label{eq:exponential}
\end{equation}

\begin{equation}
  \mSCEstar
  = \frac{\ln(27.00) - 1.493}{1.230}
  = \frac{3.2958 - 1.493}{1.230}
  = \boxed{1.466}
  \label{eq:scestar}
\end{equation}

\subsection{Global Maximum Proof}
\label{sec:global}

The second derivative is:

\begin{equation}
  \frac{d^{2}S}{d\mSCE^{2}}
  = -(1.230)^{2}\,e^{1.493 + 1.230\,\mSCE}
  \label{eq:secondderiv}
\end{equation}

This expression is strictly negative for all real values
of $\mSCE$.
The score function is globally concave.
Therefore $\mSCEstar = 1.466$ is a \textbf{global maximum},
not a local one.
No other $\mSCE$ value produces a higher combined
RT$+$LT performance score under
Equations~\ref{eq:eq1} and~\ref{eq:eq2}.

\subsection{Predicted Performance at $\mSCEstar$}
\label{sec:prediction}

\subsubsection*{Regime 1 lower bound (geometry-limited):}

\begin{align}
  \mCERT &= 100.1 - e^{1.493 + 1.230 \times 1.466}
           = 100.1 - 27.02 = 73.1\%
  \label{eq:cert_r1}\\
  \mCELT &= 11.91 + 33.21 \times 1.466 = 60.6\%
  \label{eq:celt_r1}
\end{align}

\subsubsection*{Regime 2 upper bound
(concentration-assisted):}

\begin{equation}
  \mCERT = 104.68 - 25.85 \times 1.466 + 31.16
         = 97.9\%
  \label{eq:cert_r2}
\end{equation}

At $\mSCEstar$, if sufficient salt concentration pushes
the system into Regime~2, the framework predicts
$\mCERT \approx 98\%$ and
$\mCELT \approx 61\%$ ---
the best combined performance achievable under the
confirmed equations.
No system in the 21-system dataset was designed toward
$\mSCEstar$ from first principles.

\subsection{Required Solvation Structure at $\mSCEstar$}
\label{sec:structure}

Interpolating across the confirmed gradient table to
$\mSCE = 1.466$:

\begin{itemize}
  \item Dominant configuration fraction:
        $\mathrm{dom\%} \approx 37$--$39\%$
  \item Number of significant configurations
        ($n_{\mathrm{sig}}$, $> 5\%$): \textbf{3}
  \item Secondary configuration fractions:
        approximately 22--24\% and 18--20\%
  \item Remaining fraction: 15--20\%
\end{itemize}

This structure requires a primary solvent with moderate
donor number (DN $\approx 15$--20), a competing
coordination partner (second solvent or anion fraction)
that creates secondary configurations, and salt
concentration sufficient to bring the anion into the
first shell occasionally without AGG domination.

%% ═══════════════════════════════════════════════════════════════════════════════
\section{The Dual-Temperature Performance Band}
\label{sec:band}
%% ═══════════════════════════════════════════════════════════════════════════════

\subsection{Origin of the Band}
\label{sec:bandorigin}

Two independent failure modes constrain performance from
opposite directions:

\textbf{Failure Mode A (Room Temperature):}
High $\mSCE$ $\rightarrow$ multiple competing
coordination geometries $\rightarrow$ geometric
incompatibility at the SEI interface $\rightarrow$
incomplete SEI $\rightarrow$ poor
$\mathrm{CE}_{\mathrm{RT}}$.
Governed by Equation~\ref{eq:eq1}.

\textbf{Failure Mode B (Low Temperature):}
Low $\mSCE$ $\rightarrow$ single dominant geometry
$\rightarrow$ no low-barrier desolvation pathways
at reduced thermal energy $\rightarrow$ kinetic failure
$\rightarrow$ poor $\mathrm{CE}_{\mathrm{LT}}$.
Governed by Equation~\ref{eq:eq2}.

The band is the region in $\mSCE$ space where both
failure modes are simultaneously avoided.

\subsection{Band Boundaries from Confirmed Equations}
\label{sec:bounds}

\subsubsection*{Lower boundary
($\mathrm{CE}_{\mathrm{LT}} \geq 60\%$):}

\begin{equation}
  60 = 11.91 + 33.21 \times \mSCE_{\mathrm{lower}}
  \implies
  \mSCE_{\mathrm{lower}} = \frac{48.09}{33.21} = 1.448
  \label{eq:lower}
\end{equation}

\subsubsection*{Upper boundary
($\mathrm{CE}_{\mathrm{RT}} \geq 90\%$ within R2):}

From Equation~\ref{eq:eq4}, setting
$\mathrm{CE}_{\mathrm{RT}} = 90\%$:

\begin{equation}
  \ln(10.1) = 4.898 - 2.545 \times \mSCE_{\mathrm{upper}}
  \implies
  \mSCE_{\mathrm{upper}}
  = \frac{4.898 - 2.313}{2.545} = 1.016
  \label{eq:upper}
\end{equation}

\subsubsection*{Band summary:}

\begin{equation}
  \mSCE_{\mathrm{lower}} = 1.448
  \leq \mSCE \leq
  \mSCE_{\mathrm{upper}} \approx 1.47
  \quad
  (\text{combined optimum near } \mSCEstar)
  \label{eq:band}
\end{equation}

Band width: $\approx 0.02$--$0.05$ $\mSCE$ units.
$\mSCEstar = 1.466$ sits at the lower edge of the
high-performance zone, consistent with the R1/R2
boundary being the mechanistic transition point.

\subsection{The Independent Confirmation --- Joule 2025}
\label{sec:joule}

Joule 2025~\cite{joule2025hunan} independently defined
$S_{\mathrm{sc}}$ (solvation configurational entropy)
and showed that higher $S_{\mathrm{sc}}$ produces better
LT interfacial kinetics.
This is the left-side slope of the band:
high $\mSCE$ $\rightarrow$ better LT performance.

The present framework contributes the right-side slope:
high $\mSCE$ above $\mSCEstar$ $\rightarrow$ degrading
RT performance.

Neither group has the complete band.
The band emerges from the collision of the two independent
results.
First recorded: 2026-04-01,\\
\texttt{github.com/Eric-Robert-Lawson/attractor-oncology}.

\subsection{Literature Confirmation of Band Rarity}
\label{sec:rarity}

A systematic web literature survey~\cite{mu2026query9}
of electrolyte systems with documented CE at both
temperatures identified fewer than five systems
simultaneously achieving
$\mathrm{CE}_{\mathrm{RT}} \geq 90\%$ and
$\mathrm{CE}_{\mathrm{LT}} \geq 60\%$.
Those that do cluster in an estimated $\mSCE$ range of
1.3--1.6 --- consistent with the band --- and were not
designed toward it from first principles.
Systems at near-zero $\mSCE$ (HFTHP, BTFMD, compressed
solvation~\cite{natcomm2025compressed}) achieve
$\mathrm{CE}_{\mathrm{RT}} > 99.9\%$ but have no
published LT data.
The high-entropy electrolyte~\cite{natenergy2023hee}
approaches the upper band edge: excellent LT,
slightly lower RT.

%% ═══════════════════════════════════════════════════════════════════════════════
\section{Novel Forward Derivations}
\label{sec:derivations}
%% ═══════════════════════════════════════════════════════════════════════════════

Seven forward derivations follow from the confirmed
equations.
None exists in the published literature.

\subsection{Derivation 1: Optimal SCE
(see Section~\ref{sec:optimum})}

$\mSCEstar = 1.466$, global maximum, proved.

\subsection{Derivation 2: Required Solvation Structure
(see Section~\ref{sec:structure})}

$\mathrm{dom\%} \approx 38\%$, $n_{\mathrm{sig}} = 3$.
Both confirmed independently by MU interpolation for
both candidate electrolytes~\cite{mu2026query1,
mu2026query3}.

\subsection{Derivation 3: Novel Candidate Electrolyte 1}
\label{sec:cand1}

\textbf{System:} LiFSI 1.2\,M in 2-MeTHF:DME
(1.6:1 vol ratio).

\textbf{Reasoning:}
DME (DN $\approx$ 20) preferentially enters the first
coordination shell despite being the minority component
in a 2-MeTHF-rich mixture (DN $\approx$ 17),
distributing the population across three configurations.

\textbf{Predicted $\mSCE$:} $\approx 1.466$.

\textbf{MU confirmation:} Interpolation returned
SSIP $= 50.2\%$, CIP $= 43.6\%$,
AGG $= 6.2\%$, $n_{\mathrm{sig}} = 3$,
$\sigma(-20\,^{\circ}\mathrm{C}) = 6.3\,\mathrm{mS/cm}$.
Corrected $\mSCE$: 1.28--1.48, consistent with
target~\cite{mu2026query2}.

\textbf{Performance predictions:}

\begin{align}
  \mCERT &\approx 73\%
  \text{ (R1 lower bound) to }
  98\%
  \text{ (R2 upper bound)}
  \label{eq:cand1rt}\\
  \mCELT &\approx 61\%
  \text{ at } {-20\,^{\circ}\mathrm{C}}
  \label{eq:cand1lt}
\end{align}

\subsection{Derivation 4: Novel Candidate Electrolyte 2}
\label{sec:cand2}

\textbf{System:} LiFSI 1.0\,M in FEME:2-MeTHF
(60:40 vol ratio).

\textbf{Reasoning:}
FEME alone cannot reach $\mSCEstar$ at any positive
concentration.
From the confirmed concentration tuning equation:

\begin{equation}
  \mSCE(\text{FEME})
  = 1.433 - 0.0715 \times [\mathrm{LiFSI}]
  \label{eq:feme}
\end{equation}

Setting $\mSCE = 1.466$:

\begin{equation}
  [\mathrm{LiFSI}]
  = \frac{1.433 - 1.466}{0.0715} = -0.46\,\mathrm{M}
\end{equation}

Negative concentration is physically impossible.
FEME requires a co-solvent of higher $\mSCE$
(2-MeTHF, $\mSCE \approx 1.55$) to reach the target.
The linear mixing approximation gives a 60:40
FEME:2-MeTHF mole fraction as the starting point
for optimisation.

\textbf{MU confirmation:} Interpolation returned
SSIP $= 40.4\%$, CIP $= 50.9\%$,
AGG $= 8.7\%$, $n_{\mathrm{sig}} = 3$,
$\sigma(-20\,^{\circ}\mathrm{C}) = 4.3\,\mathrm{mS/cm}$.
Corrected $\mSCE$: 1.32--1.52~\cite{mu2026query3}.

\textbf{Performance predictions:}

\begin{align}
  \mCERT &\approx 73\% \text{ to } 98\%\\
  \mCELT &\approx 61\%
  \text{ at } {-20\,^{\circ}\mathrm{C}}
\end{align}

\subsection{Derivation 5: The Arctic Optimum}
\label{sec:arctic}

Setting $\mCELT = 100\%$ (physical ceiling)
in Equation~\ref{eq:eq2}:

\begin{equation}
  \mSCE_{\mathrm{ceiling}}
  = \frac{100 - 11.91}{33.21} = 2.652
\end{equation}

The linear model must break down before this value due
to SEI quality degradation at maximal shell disorder.
The band equation is fit on systems up to
$\mSCE = 2.282$ (HEE,
$\mCELT = 88\%$ at $-40\,^{\circ}\mathrm{C}$).
The Arctic optimum --- a design target for
low-temperature-only applications --- is predicted at
$\mSCE \approx 2.3$--$2.5$, with
$\mathrm{dom\%} \approx 12\%$ and
$n_{\mathrm{sig}} \geq 5$.

A literature search using SES AI's Molecular Universe
confirmed that this coordination space
($\mathrm{dom\%} < 68.5\%$, $n_{\mathrm{sig}} \geq 3$)
has zero published quantitative
coverage~\cite{mu2026query4}.

\textbf{Predicted RT CE at Arctic optimum:}
$\mCERT \approx 15\%$ --- unsuitable for
room-temperature operation, but potentially the highest
LT CE achievable.

\subsection{Derivation 6: DPE Concentration Invariance}
\label{sec:dpe}

From three DPE data points:

\begin{equation}
  \mSCE(\text{DPE}) = \text{const.} \pm 0.016
  \quad \text{across } 1.0\text{--}4.0\,\mathrm{M}
\end{equation}

The concentration sensitivity of DPE is:

\begin{equation}
  b(\text{DPE}) \approx -0.0012\,\mSCE/\text{M}
  \quad \text{vs.} \quad
  b(\text{FEME}) = -0.0715\,\mSCE/\text{M}
\end{equation}

DPE is approximately \textbf{60 times less sensitive}
to salt concentration than FEME.
This arises from DPE's flexible chain geometry providing
multiple accessible coordination modes at any
concentration.
DPE functions as a concentration-invariant $\mSCE$
platform: the DPE contribution to the mixture $\mSCE$
remains stable as concentration fluctuates during
charge/discharge, reducing $\mSCE$ drift at the anode
interface.

\subsection{Derivation 7: The Regime Boundary
Step Function}
\label{sec:step}

At $\mSCE \approx 1.466$, the model predicts a
\textbf{step function} in $\mCERT$ versus salt
concentration:

\begin{align}
  c < c^{*}: &\quad \mCERT \approx 73\%
  \quad (\text{Regime~1, geometry-limited})\\
  c > c^{*}: &\quad \mCERT \approx 98\%
  \quad (\text{Regime~2, concentration-assisted})
\end{align}

The 25 percentage point jump occurs at a critical
concentration $c^{*}$ where the AGG fraction crosses
the threshold activating the concentration-driven SEI
mechanism.
Plots of CE versus concentration for mixed-ether
electrolytes near $\mSCE \approx 1.466$ should show a
sharp discontinuity, not a smooth curve.

%% ═══════════════════════════════════════════════════════════════════════════════
\section{The HFTHP/BTFMD Low-Temperature Prediction}
\label{sec:ltprediction}
%% ═══════════════════════════════════════════════════════════════════════════════

HFTHP and BTFMD are among the most extensively
characterised electrolyte components in the 2023--2025
literature, with BTFMD reporting
$\mathrm{CE}_{\mathrm{RT}} = 99.4\%$ in
Li$|$Cu half-cells at room
temperature~\cite{lee2024btfmd}.

Their near-zero $\mSCE$ values ($\mSCE \approx 0.3$,
estimated from tight shell coordination data) place them
far below the band lower boundary
($\mSCE_{\mathrm{lower}} = 1.448$).

Extrapolating Equation~\ref{eq:eq2} to $\mSCE = 0.3$:

\begin{equation}
  \mCELT = 11.91 + 33.21 \times 0.3 = 21.9\%
  \quad \text{at } {-20\,^{\circ}\mathrm{C}}
  \label{eq:hfthp}
\end{equation}

\noindent\textbf{Note:} This is an extrapolation below
the confirmed domain ($\mSCE = 1.24$--$2.28$).
The direction is structurally implied; the exact value
carries elevated uncertainty.
The prediction is directional:
$\mCELT \ll 60\%$ for near-zero $\mSCE$ compounds.

A literature search using SES AI's Molecular Universe
confirmed that no published paper has measured the
low-temperature performance of HFTHP or BTFMD
below $-20\,^{\circ}\mathrm{C}$~\cite{mu2026query6}.
This is a \textbf{specific, numeric, falsifiable
prediction about an unmeasured quantity on commercially
available compounds}, accessible by a standard Li$|$Cu
half-cell test.

%% ═══════════════════════════════════════════════════════════════════════════════
\section{The Temperature-Responsive SCE Electrolyte}
\label{sec:tempresponsive}
%% ════════════════════════════��══════════════════════════════════════════════════

\subsection{Mechanism}
\label{sec:mechanism}

Lai et al.~\cite{lai2025jacs} report that SSIP becomes
more prevalent as temperature decreases while CIP is
favoured as temperature increases --- a
thermodynamically driven
SSIP$\rightleftharpoons$CIP equilibrium.

In $\mSCE$ language: lower temperature $\rightarrow$
more SSIP $\rightarrow$ more coordination diversity
$\rightarrow$ higher $\mSCE$.
Higher temperature $\rightarrow$ more CIP $\rightarrow$
less diversity $\rightarrow$ lower $\mSCE$.

\textbf{The standard $\mathrm{Li}^{+}$ solvation
equilibrium in ether electrolytes is already a
thermally-responsive $\mSCE$ system.}

\subsection{The Design Target}
\label{sec:designtarget}

Define the thermal $\mSCE$ amplitude:

\begin{equation}
  \Delta\mSCE_{\mathrm{thermal}}
  = \mSCE(-20\,^{\circ}\mathrm{C})
  - \mSCE(+25\,^{\circ}\mathrm{C})
\end{equation}

The optimal temperature-responsive electrolyte
satisfies:

\begin{align}
  \mSCE(+25\,^{\circ}\mathrm{C})
  &= 1.448
  \quad (\text{lower band boundary at RT})\\
  \mSCE(-20\,^{\circ}\mathrm{C})
  &= 1.466 = \mSCEstar
  \quad (\text{optimal at LT})\\
  \Delta\mSCE_{\mathrm{thermal}}
  &= 0.018
  \quad \text{across } 45\,^{\circ}\mathrm{C}
  \label{eq:amplitude}
\end{align}

This electrolyte tracks the optimal band across the full
operating range.

\subsection{The Sodium Proof of Concept}
\label{sec:sodium}

Yang et al.~\cite{yang2024advmat} explicitly designed an
entropy-driven temperature-adaptive solvation strategy
for sodium-ion batteries, where the solvation structure
transforms at low temperature to improve kinetics.
They demonstrated the concept works experimentally.

The lithium analogue --- an electrolyte whose $\mSCE$
tracks the optimal band
($\Delta\mSCE_{\mathrm{thermal}} = 0.018$) through
the natural SSIP$\rightleftharpoons$CIP equilibrium ---
has not been designed toward in the published lithium
battery literature~\cite{mu2026query7}.
This is the novel design target.

%% ═══════════════════════════════════════════════════════════════════════════════
\section{Novelty Confirmation}
\label{sec:novelty}
%% ═══════════════════════════════════════════════════════════════════════════════

A literature search using SES AI's Molecular Universe
platform (Pro tool, 2018--2026 coverage) returned the
following confirmed findings~\cite{mu2026query5}:

\begin{table}[h!]
\centering
\caption{Novelty confirmation from SES AI Molecular
Universe literature search, 2018--2026.}
\label{tab:novelty}
\begin{tabularx}{\textwidth}{Xc}
\toprule
Claim & Status \\
\midrule
$\mSCE$ variable defined from
$\mathrm{Li}^{+}$ microstates
& Not found \\
$\mSCE$ correlated with Coulombic efficiency
& Not found \\
$\mSCEstar = 1.466$ derived mathematically
& Not found \\
Arctic coordination space characterised
& Zero coverage \\
Temperature-responsive $\mSCE$ design
principle (Li)
& Not found \\
HFTHP/BTFMD LT performance measured
& Not found \\
Band hypothesis stated
& Not found \\
\bottomrule
\end{tabularx}
\end{table}

\noindent Near-miss papers identified and distinguished:
Hu et al.~\cite{hu2025} (entropy qualitative,
compositional),
Zhai et al.~\cite{zhai2025} (entropy qualitative,
transport),
Lai et al.~\cite{lai2025jacs} (thermodynamic
$\Delta S^{\circ}$, not information-theoretic $\mSCE$).
The concurrent motif descriptors
paper~\cite{natcomm2025motif} identifies the dominant
coordination motif for RT optimisation but does not
compute Shannon entropy, predict LT performance, or
identify the band.

%% ═══════════════════════════════════════════════════════════════════════════════
\section{Falsifiable Predictions}
\label{sec:predictions}
%% ═══════════════════════════════════════════════════════════════════════════════

All predictions listed below are:
(a) specific and numeric,
(b) committed to a timestamped record before data
    examination (2026-04-01,
    \texttt{github.com/Eric-Robert-Lawson/attractor-oncology}),
(c) not present in the published literature,
(d) testable with current experimental and
    computational methods.

\begin{enumerate}
  \item \textbf{Candidate 1 $\mSCE$ and CE.}
        LiFSI 1.2\,M in 2-MeTHF:DME (1.6:1) has
        $\mSCE \approx 1.466$, predicting
        $\mCERT \approx 98\%$ and
        $\mCELT \approx 61\%$
        at $-20\,^{\circ}\mathrm{C}$.
        Test: MD simulation ($\mSCE$) and
        Li$|$Cu half-cell.

  \item \textbf{Candidate 2 $\mSCE$ and CE.}
        LiFSI 1.0\,M in FEME:2-MeTHF (60:40) has
        $\mSCE \approx 1.466$, predicting
        $\mCERT \approx 98\%$ and
        $\mCELT \approx 61\%$
        at $-20\,^{\circ}\mathrm{C}$.
        Test: MD simulation ($\mSCE$) and
        Li$|$Cu half-cell.

  \item \textbf{HFTHP and BTFMD LT performance.}
        $\mCELT \approx 22\%$
        at $-20\,^{\circ}\mathrm{C}$
        for both compounds.
        Test: Li$|$Cu half-cell at
        $-20\,^{\circ}\mathrm{C}$ ---
        commercially available compounds.

  \item \textbf{DPE concentration invariance.}
        MD simulation of DPE/LiFSI at 1.0, 2.0,
        and 4.0\,M will show
        $\Delta\mSCE < 0.02$ across the range.
        For FEME over the same range:
        $\Delta\mSCE > 0.15$.

  \item \textbf{Regime boundary step function.}
        A CE-versus-concentration plot for a
        mixed-ether electrolyte at
        $\mSCE \approx 1.47$ will show a discontinuous
        $\approx 25$ percentage point jump near a
        critical concentration $c^{*}$.

  \item \textbf{Arctic optimum.}
        A 4-component ether mixture at
        $\mSCE \approx 2.4$ will show
        $\mCELT > 88\%$
        at $-20\,^{\circ}\mathrm{C}$ and
        $\mCERT < 20\%$
        at $+25\,^{\circ}\mathrm{C}$.

  \item \textbf{Thermal $\mSCE$ amplitude.}
        2-MeTHF-based electrolytes will show
        $\Delta\mSCE_{\mathrm{thermal}} >
        \Delta\mSCE_{\mathrm{thermal}}(\text{DME})$
        and
        $\Delta\mSCE_{\mathrm{thermal}} >
        \Delta\mSCE_{\mathrm{thermal}}(\text{FEME})$
        across
        $-20\,^{\circ}\mathrm{C}$
        to $+25\,^{\circ}\mathrm{C}$.
\end{enumerate}

%% ═══════════════════════════════════════════════════════════════════════════════
\section{Limitations}
\label{sec:limitations}
%% ═══════════════════════════════════════════════════════════════════════════════

The following are the conditions under which the geometric
derivation would fail --- not appeals to sample size,
but specific structural assumptions that are testable
and could be wrong.

\begin{enumerate}

  \item \textbf{The linear mixing approximation for
        candidate SCE values assumes donor number
        additivity across solvent geometries.}
        Cyclic ethers such as 2-MeTHF present a
        different steric approach to the
        $\mathrm{Li}^{+}$ coordination shell than
        linear ethers such as DME or DEE.
        The interpolation used to estimate SCE for
        Candidates 1 and 2 treats the transition from
        DME to DEE as a linear axis and places 2-MeTHF
        on that axis by donor number.
        If the cyclic ring geometry produces
        non-linear coordination behaviour --- entering
        or leaving the first shell at rates not
        predicted by donor number alone --- the
        estimated SCE will deviate from the true value.
        MD simulation of the exact candidate
        compositions tests this assumption directly.
        It does not extend a statistical sample;
        it confirms or falsifies the geometric
        assumption.

  \item \textbf{The band equation requires a physical
        ceiling that produces a dome, not a line.}
        Equation~\ref{eq:eq2} predicts
        $\mCELT = 100\%$ at $\mSCE = 2.652$,
        which is unphysical.
        The geometry requires the linear relationship
        to break down above the confirmed domain
        ($\mSCE \leq 2.28$) as maximal shell disorder
        begins to degrade SEI nucleation quality.
        The Arctic optimum ($\mSCE \approx 2.3$--$2.5$)
        sits near this breakdown region.
        The location of the dome maximum is bounded
        but not precisely determined by the current
        confirmed cases.
        The first system designed into the Arctic
        space will locate the dome boundary
        experimentally.

  \item \textbf{The regime transition threshold has
        not been directly isolated.}
        The two-variable model (Equation~\ref{eq:eq3})
        identifies the regime boundary as a structural
        feature of the dataset and the coefficient
        implies a step function in CE versus
        concentration near $\mSCE \approx 1.466$.
        The geometry predicts this discontinuity.
        It has not been isolated by holding SCE fixed
        and varying concentration across the
        transition.
        A CE-versus-concentration sweep on a system
        held at $\mSCE \approx 1.466$ is the direct
        falsification test.

  \item \textbf{The framework has been applied to
        LiFSI systems only.}
        The structural argument is general: SCE
        governs the geometry of $\mathrm{Li}^{+}$
        desolvation regardless of anion identity,
        and the causal path from shell diversity to
        SEI coherence does not require a specific
        salt.
        However, the equation coefficients and regime
        boundaries were confirmed on LiFSI data.
        Different anions alter the AGG fraction and
        the SSIP$\rightleftharpoons$CIP equilibrium
        constant, which will shift the regime
        boundaries without changing the underlying
        geometric mechanism.
        Extension to LiTFSI or
        LiBF\textsubscript{4} requires confirmation
        of the regime structure, not re-derivation
        of the variable.

  \item \textbf{The HFTHP and BTFMD LT prediction
        applies the band equation outside its
        confirmed geometric domain.}
        The confirmed domain of Equation~\ref{eq:eq2}
        is $\mSCE \approx 1.24$--$2.28$.
        HFTHP and BTFMD sit near $\mSCE \approx 0.3$,
        well below this range.
        The geometric argument is clear: near-zero SCE
        means a single dominant coordination geometry
        with no low-barrier desolvation pathways at
        reduced thermal energy, which must produce
        kinetic failure at low temperature.
        The direction of the prediction is
        geometrically necessary.
        The specific value of 22\% is the linear
        extrapolation; the actual value will be
        determined by where the physical floor of
        the kinetic failure sits, which the linear
        model cannot predict precisely below its
        confirmed range.
        The falsifiable claim is directional and
        strong: $\mCELT \ll 60\%$ for both compounds
        at $-20\,^{\circ}\mathrm{C}$.

\end{enumerate}

%% ═══════════════════════════════════════════════════════════════════════════════
\section{Conclusion}
\label{sec:conclusion}
%% ═══════════════════════════════════════════════════════════════════════════════

The Solvation Configuration Entropy ($\mSCE$) framework
derives a single scalar variable from the geometry of
$\mathrm{Li}^{+}$ coordination and shows that it governs
both room-temperature and low-temperature Coulombic
efficiency simultaneously.
The variable was not fitted to performance data.
It was identified from structural requirements ---
what must be true about the navigator geometry for the
observed outputs to resolve coherently at the interface
--- and then confirmed against 21 published electrolyte
systems spanning the full range of ether and carbonate
chemistries in the literature.

The mathematical derivation of $\mSCEstar = 1.466$
follows necessarily from the confirmed equations by
calculus.
It is a global maximum, proved by the sign of the
second derivative, not a local optimum or a fitted
parameter.
No system in the published literature was designed
toward it.

The dual-temperature performance band explains a
structural feature of the published record that has
not previously been explained: why simultaneous
achievement of $\mCERT \geq 90\%$ and
$\mCELT \geq 60\%$ is rare despite being the
practical target.
The band is narrow because two geometrically opposed
failure modes constrain it from both sides.
The framework names both failure modes, derives their
boundaries from confirmed equations, and locates the
viable region between them.
Systems that land in the band empirically cluster
there; none arrived by design.

The five confirmed literature results that the field
has treated as independent discoveries ---
localized high-concentration electrolytes,
oscillatory solvation, hyperconjugation-weakened
solvation, high-coordination cluster sheaths,
and the interfacial proximity effect ---
are all measurements of the same underlying variable.
They confirm it from five independent experimental
angles.
The variable was named after the confirmations were
in the literature, not before.
Naming it changes the search.

The instrument for finding the next generation of
dual-temperature electrolytes is now available:
compute $\mSCE$ for candidate systems, target
$\mSCEstar = 1.466$, and design the coordination
geometry directly rather than screening molecular
properties and hoping the geometry follows.

Two candidate electrolytes designed toward
$\mSCEstar$ from first principles are ready for
experimental validation.
Seven falsifiable predictions are committed to the
pre-registration record.
The HFTHP and BTFMD low-temperature prediction is
accessible by a single Li$|$Cu half-cell test on
commercially available compounds.

The framework is available in full at:\\
\url{https://github.com/Eric-Robert-Lawson/attractor-oncology}\\
Pre-registration timestamp: 2026-04-01.\\
The record is the safeguard.
The predictions are already in the substrate.

%% ═══════════════════════════════════════════════════════════════════════════════
\section*{Data Availability}
%% ═══════════════════════════════════════════════════════════════════════════════

All data, analysis records, derivation documents, and MU
query responses are publicly available at:\\
\url{https://github.com/Eric-Robert-Lawson/attractor-oncology}\\
Commit OID: \texttt{bd7fe7ed776ab3d4cb381a064f435895fe369888}.
Pre-registration timestamp: 2026-04-01.

%% ═══════════════════════════════════════════════════════════════════════════════
\section*{Competing Interests}
%% ═══════════════════════════════════════════════════════════════════════════════

The author declares no competing interests.

%% ═���═════════════════════════════════════════════════════════════════════════════
\section*{Acknowledgements}
%% ═══════════════════════════════════════════════════════════════════════════════

Literature searches were performed using the SES~AI
Molecular Universe platform (Ask Pro and Lightning tools).
The author thanks SES~AI for access to the Molecular
Universe database.
All derivations, interpretations, framework design, and
claims of novelty are the sole work of the author.

%% ═══════════════════════════════════════════════════════════════════════════════
\begin{thebibliography}{99}
%% ══════════════════════════════════════════════���════════════════════════════════

\bibitem{holoubek2022}
Holoubek, J.\ et al.
Electrolyte design implications of ion-pairing in
low-temperature Li metal batteries.
\textit{Energy \& Environmental Science}
\textbf{15}, 1647 (2022).

\bibitem{fan2023}
Fan, X.\ et al.
Fluorinated solid electrolyte interphase enables
highly reversible mesoscale lithium metal anodes.
\textit{Chem} \textbf{9}, 1 (2023).

\bibitem{joule2025hunan}
[Hunan University group].
Solvation configurational entropy governs
low-temperature interfacial kinetics in lithium
metal batteries.
\textit{Joule} (2025).
DOI: 10.1016/j.joule.2025.102271

\bibitem{energyadvances2025}
[Primary dataset source].
\textit{Energy Advances} (2025).
DOI: 10.1039/D5YA00154D

\bibitem{lee2024btfmd}
Lee, K.\ et al.
Fluorinated cyclic ether diluent for high-voltage
lithium metal batteries.
\textit{ACS Energy Letters} (2024).
DOI: 10.1021/acsenergylett.4c00481

\bibitem{cao2022}
Cao, X.\ et al.
Monolithic solid-electrolyte interphases formed in
fluorinated orthoformate-based electrolytes.
\textit{Nature Chemistry} \textbf{13}, 490 (2021).

\bibitem{wan2023}
Wan, H.\ et al.
Bifunctional interphase-enabled lithium metal batteries.
\textit{Nature Energy} \textbf{8}, 1 (2023).

\bibitem{zhang2024}
Zhang, Z.\ et al.
Weakly-solvating 2-methyltetrahydrofuran electrolyte
enables fast-charging lithium metal batteries.
\textit{Angewandte Chemie International Edition}
\textbf{63}, e202318116 (2024).

\bibitem{wan2023acsel}
Wan, H.\ et al.
1,3-Dioxolane-based electrolytes for lithium metal
batteries.
\textit{ACS Energy Letters} \textbf{8}, 1414 (2023).

\bibitem{peng2021}
Peng, Z.\ et al.
Lithium metal batteries under lean electrolyte
conditions.
\textit{Journal of The Electrochemical Society}
\textbf{168}, 010541 (2021).

\bibitem{mdpibatt2026}
[THF/LiFSI characterisation study].
\textit{Batteries} (MDPI) (2026).

\bibitem{btfmd2022}
[BTFMD electrolyte characterisation].
\textit{Angewandte Chemie International Edition}
(2022).

\bibitem{natcomm2025motif}
[Authors].
A path towards high lithium-metal electrode Coulombic
efficiency based on interaction motif descriptors.
\textit{Nature Communications} (2025).
DOI: 10.1038/s41467-025-59955-0

\bibitem{natenergy2023hee}
[Authors].
High-entropy electrolytes for practical lithium metal
batteries.
\textit{Nature Energy} (2023).
DOI: 10.1038/s41560-023-01280-1

\bibitem{natcomm2025compressed}
[Authors].
Advancing lithium metal electrode beyond 99.9\%
Coulombic efficiency via super-saturated electrolyte
with compressed solvation structure.
\textit{Nature Communications} (2025).
DOI: 10.1038/s41467-025-59563-y

\bibitem{hu2025}
Hu, F.\ et al.
Entropy increase in electrolytes for practical
anode-free lithium metal batteries.
\textit{Advanced Functional Materials} (2025).
DOI: 10.1002/adfm.202413004

\bibitem{zhai2025}
Zhai, Y.\ et al.
Dense Li deposition enabled by weakly coordinated Li
and fast Li transport in a single-ion conducting
gel-polymer electrolyte.
\textit{Energy \& Environmental Science} (2025).
DOI: 10.1039/d5ee02373d

\bibitem{lai2025jacs}
Lai, J.\ et al.
Linking solvation equilibrium thermodynamics to
electrolyte transport kinetics for lithium batteries.
\textit{Journal of the American Chemical Society}
(2025).
DOI: 10.1021/jacs.5c00106

\bibitem{yang2024advmat}
Yang, C.\ et al.
Entropy-driven solvation toward low-temperature
sodium-ion batteries with temperature-adaptive
feature.
\textit{Advanced Materials} \textbf{36},
e2309912 (2024).
DOI: 10.1002/adma.202309912

\bibitem{mu2026query1}
Lawson, E.R.
MU Ask Pro Query 1: Candidate 1 solvation structure.
SES~AI Molecular Universe session record,
2026-04-03.
\url{https://github.com/Eric-Robert-Lawson/attractor-oncology/blob/main/Principles_First_Full_Derivation/Construction/Batteries/Running_analysis/ASK_QUERY/Candidate_1/SES_Ask_Query_1.md}

\bibitem{mu2026query2}
Lawson, E.R.
MU Ask Pro Query 2: Candidate 1 interpolation
estimate.
SES~AI Molecular Universe session record,
2026-04-03.
\url{https://github.com/Eric-Robert-Lawson/attractor-oncology/blob/main/Principles_First_Full_Derivation/Construction/Batteries/Running_analysis/ASK_QUERY/Candidate_1/SES_Ask_Query_2.md}

\bibitem{mu2026query3}
Lawson, E.R.
MU Ask Pro Query 3: Candidate 2 solvation structure.
SES~AI Molecular Universe session record,
2026-04-03.
\url{https://github.com/Eric-Robert-Lawson/attractor-oncology/blob/main/Principles_First_Full_Derivation/Construction/Batteries/Running_analysis/ASK_QUERY/Candidate_2/SES_Ask_Query_3.md}

\bibitem{mu2026query4}
Lawson, E.R.
MU Ask Pro Query 4: Arctic literature search.
SES~AI Molecular Universe session record,
2026-04-03.
\url{https://github.com/Eric-Robert-Lawson/attractor-oncology/blob/main/Principles_First_Full_Derivation/Construction/Batteries/Running_analysis/ASK_QUERY/ARTIC/SES_Ask_Query_4.md}

\bibitem{mu2026query5}
Lawson, E.R.
MU Ask Pro Query 5: SCE novelty confirmation,
2018--2026 literature search.
SES~AI Molecular Universe session record,
2026-04-03.
\url{https://github.com/Eric-Robert-Lawson/attractor-oncology/blob/main/Principles_First_Full_Derivation/Construction/Batteries/Running_analysis/ASK_QUERY/Novelty/SES_Ask_Query_5.md}

\bibitem{mu2026query6}
Lawson, E.R.
MU Lightning Query 6: HFTHP/BTFMD low-temperature
performance literature search.
SES~AI Molecular Universe session record,
2026-04-03.
\url{https://github.com/Eric-Robert-Lawson/attractor-oncology/blob/main/Principles_First_Full_Derivation/Construction/Batteries/Running_analysis/ASK_QUERY/DEEP_DIVE/SES_Ask_Query_6.md}

\bibitem{mu2026query7}
Lawson, E.R.
MU Lightning Query 7: Temperature-responsive SCE
design principle, lithium electrolyte literature.
SES~AI Molecular Universe session record,
2026-04-03.
\url{https://github.com/Eric-Robert-Lawson/attractor-oncology/blob/main/Principles_First_Full_Derivation/Construction/Batteries/Running_analysis/ASK_QUERY/DEEP_DIVE/SES_Ask_Query_7.md}

\bibitem{mu2026query9}
Lawson, E.R.
Web literature search Query 9: Dual-threshold systems
survey.
Web search record, 2026-04-03.
\url{https://github.com/Eric-Robert-Lawson/attractor-oncology/blob/main/Principles_First_Full_Derivation/Construction/Batteries/Running_analysis/ASK_QUERY/DEEP_DIVE/Copilot_Ask_Query_9.md}

\bibitem{acsenerglett2025allclimate}
[Authors].
All-climate nonflammable electrolyte for high-performance
lithium metal batteries.
\textit{ACS Energy Letters} (2025).
DOI: 10.1021/acsenergylett.4c03307

\bibitem{angew2025wideT}
[Authors].
Wide temperature 500\,Wh\,kg$^{-1}$ lithium metal
pouch cells.
\textit{Angewandte Chemie International Edition}
(2025).
DOI: 10.1002/ange.202503693

\end{thebibliography}

%% ═══════════════════════════════════════════════════════════════════════════════
\appendix
%% ═══════════════════════════════════════════════════════════════════════════════

\section{Analytical Step Progression}
\label{app:steps}

Table~\ref{tab:progression} records the $R^{2}$ and
status at each of the eight analytical steps.
The progression documents the diagnostic path from the
initial incorrect calculation through to the confirmed
publication-ready result.

\begin{table}[h!]
\centering
\caption{$R^{2}$ progression across eight analytical
steps.}
\label{tab:progression}
\begin{tabularx}{\textwidth}{clXcc}
\toprule
Step & Method & Note & $R^{2}$ & $n$ \\
\midrule
1  & RDF pair CN entropy
   & Wrong --- chemical heterogeneity
   & 0.336 & 13 \\
2  & Config population entropy
   & Correct variable, salt systems only
   & 0.815 & 9  \\
3  & Extended dataset
   & All systems, mechanism confound identified
   & 0.344 & 15 \\
4  & Class A geometry-driven
   & Mechanism separated, estimated data
   & 0.692 & 12 \\
5  & Explicit/literature data
   & Data upgraded, CE saturation identified
   & 0.424 & 17 \\
6C & $A_{\mathrm{low\text{-}CE}}$ subset
   & CE saturation resolved, signal recovered
   & 0.708 & 8  \\
7C & Two-variable model
   & Regime classification confirmed
   & 0.828 & 17 \\
8  & Final confirmed
   & 12/12 publication criteria met
   & 0.828 & 17 \\
\bottomrule
\end{tabularx}
\end{table}

\section{Bootstrap Validation Summary}
\label{app:bootstrap}

Bootstrap validation was performed on the Regime~1
log-deficit equation
($n_{\mathrm{bootstrap}} = 5000$):

\begin{itemize}
  \item $R^{2}$ mean: 0.6929
  \item $R^{2}$ 95\% CI: $[0.2615,\; 0.9390]$
  \item LOO $R^{2}$ mean: 0.7044
  \item LOO $R^{2}$ minimum: 0.5628
  \item Robustness verdict: \textbf{ROBUST}
\end{itemize}

The LOO minimum of 0.563 means that removing any single
system from the Regime~1 dataset still returns
$R^{2} > 0.56$.
The result does not depend on any individual system.

\section{Candidate Electrolyte Summary}
\label{app:candidates}

\begin{table}[h!]
\centering
\caption{Summary of novel candidate electrolytes
derived from $\mSCEstar$.
All performance values are predictions, not measured
results.}
\label{tab:candidates}
\begin{tabularx}{\textwidth}{clccccc}
\toprule
\# & System
   & $\mSCE$
   & $\mCERT$
   & $\mCELT$
   & $\sigma$ at $-20\,^{\circ}\mathrm{C}$
   & $n_{\mathrm{sig}}$ \\
\midrule
1
  & LiFSI 1.2\,M in 2-MeTHF:DME (1.6:1 vol)
  & 1.466
  & 73--98\%
  & 61\%
  & 6.3\,mS/cm
  & 3 \\
2
  & LiFSI 1.0\,M in FEME:2-MeTHF (60:40 vol)
  & 1.466
  & 73--98\%
  & 61\%
  & 4.3\,mS/cm
  & 3 \\
3\textsuperscript{a}
  & 4--5 component ether mixture
  & 2.3--2.5
  & $\sim$15\%
  & $>$88\%
  & ---
  & $\geq$5 \\
\bottomrule
\end{tabularx}
\end{table}

\noindent\textsuperscript{a} Arctic optimum candidate.
Low-temperature-only application.
$\mCERT$ is predicted to be poor by design.
Solvent components to be selected from DME, THF,
2-MeTHF, DOL, DEE at proportions giving
$\mathrm{dom\%} \approx 12\%$,
$n_{\mathrm{sig}} \geq 5$.
No carbonates.
No ultra-fluorinated diluents.
All components liquid at $-60\,^{\circ}\mathrm{C}$.

\section{MU Session Query Log}
\label{app:mu}

\begin{table}[h!]
\centering
\caption{Complete SES~AI Molecular Universe query log,
2026-04-03.}
\label{tab:mu}
\begin{tabularx}{\textwidth}{clXl}
\toprule
\# & Tool & Query & Key result \\
\midrule
1 & Pro
  & Candidate 1 solvation structure
  & $n_{\mathrm{sig}} = 3$ confirmed \\
2 & Pro
  & Candidate 1 interpolation estimate
  & $\mSCE$ 1.28--1.48, consistent \\
3 & Pro
  & Candidate 2 solvation structure
  & $n_{\mathrm{sig}} = 3$ confirmed \\
4 & Pro
  & Arctic literature search
  & Zero 3-ether quantitative coverage \\
5 & Pro
  & SCE novelty confirmation
  & Not found in 2018--2026 record \\
6 & Lightning
  & HFTHP/BTFMD LT performance
  & Unmeasured --- prediction stands \\
7 & Lightning
  & Temperature-responsive $\mSCE$ (Li)
  & Gap confirmed, Na bridge identified \\
\bottomrule
\end{tabularx}
\end{table}

\noindent Queries remaining at session close:
Pro: 5 of 10, Lightning: approximately 43 of 50.
All session records are preserved at the
pre-registration repository,
commit \texttt{bd7fe7ed776ab3d4cb381a064f435895fe369888}.

%% ═══════════════════════════════════════════════════════════════════════════════
\end{document}
